\documentclass[11pt]{article}
\usepackage[T1]{fontenc}
\usepackage{lmodern}
\usepackage{amsmath,amssymb}
\usepackage{booktabs}
\usepackage{array}
\usepackage[a4paper,margin=2.2cm]{geometry}
\usepackage[hidelinks]{hyperref}
\newcommand{\doi}[1]{\href{https://doi.org/#1}{\texttt{#1}}}
\setlength{\parskip}{3pt}
\title{\vspace{-1.5cm}\textbf{Packages v4 --- what was delivered, and what the searches changed}\\[2pt]
\large Milestone report}
\author{Dr.\ Vipin \\ Gurugram University}
\date{30 October 2026}
\begin{document}
\maketitle
\vspace{-8mm}
\footnotesize

\section*{1. What is in the delivery}

Three archives, each with exactly two \LaTeX\ files and nothing else.

\begin{center}\footnotesize
\begin{tabular}{lll}
\toprule
archive & file 1 & file 2\\
\midrule
\texttt{Lamb\_Sumit.zip} & \texttt{01\_Research\_Plan\_Sumit.tex} & \texttt{02\_Basics\_and\_Definitions\_Sumit.tex}\\
\texttt{Lamb\_Neha.zip} & \texttt{01\_Research\_Plan\_Neha.tex} & \texttt{02\_Basics\_and\_Definitions\_Neha.tex}\\
\texttt{Lamb\_Anu.zip} & \texttt{01\_Research\_Plan\_Anu.tex} & \texttt{02\_Basics\_and\_Definitions\_Anu.tex}\\
\bottomrule
\end{tabular}
\end{center}

Each plan carries the seven sections you asked for: title, abstract and keywords; introduction with cited
literature and the gap; problem and assumptions; the model up to but not including the boundary conditions;
a six-step work plan with the switching-off tolerance written into it; figures, risks and timeline; and
references. The notation table sits inside the model section and carries units. Each basics document is
written for a reader at class ten, has no equations to derive and no numbers, and ends in a glossary with
units. No LaTeX engine exists in this workspace, so nothing is compiled; the code is kept to the eight
packages you allowed, with no images and no custom environments.

\section*{2. The seven corrections, and where each one landed}

\begin{center}\footnotesize
\begin{tabular}{p{0.4cm}p{4.6cm}p{8.6cm}}
\toprule
\# & correction & what was done\\
\midrule
1 & NSGT in Lim form, momentum form, equation number & (N1) is $(1-l_{\rm nl}^2\nabla^2)t_{ij}=\sigma_{ij}$ with the gradient operator on $c{:}\varepsilon$ only. (N2) is the equivalent momentum form with the operator on the inertia. Equations (18) and (19) of Lim et al.\ are cited \emph{through} Barretta \& Marotti de Sciarra (2018), because the Elsevier file was closed to me; the status is [S], secondary, and the scholar is told to confirm (18) on the page.\\
2 & local couplings & the gradient operator removed from A2, B3, C2, C3. It survives only on $c{:}\varepsilon$. The modelling choice is stated once, with the three reasons: $D$ and $B$ stay second order, the extra face conditions stay mechanical, and no field gradient is differentiated twice.\\
3 & A6 Peltier term, units, sign, recombination & the band-gap term is gone. (A5) carries $\beta_{\rm th}\nabla^2\theta$ with $\beta_{\rm th}=b_nn_0S_n$, unit m$^{-1}$s$^{-1}$K$^{-1}$, negative for electrons. (A6) carries $-q\Pi_{\rm th}\partial_t\Delta n$ with $\Pi_{\rm th}=T_0S_n$ and (A7) states the Onsager ratio. Recombination is an optional pair that vanishes as $\tau_R\to\infty$. Every term of (A5) and (A6) is dimension-tested in a table.\\
4 & three claims removed or sourced & all three removed; the searches that replaced them are in section 3 below.\\
5 & Lubensky and phason dynamics & corrected to Lubensky, Ramaswamy \& Toner, \emph{Phys.\ Rev.\ B} \textbf{32} (1985) 7444, icosahedral, cited for the mechanism only. Added Agiasofitou \& Lazar, \emph{Int.\ J.\ Solids Struct.} \textbf{51} (2014) 923, Eqs.~(26) and (30), for the telegraph phason balance, and Baggioli \& Landry, \emph{SciPost Phys.} \textbf{9} (2020) 062, Eq.~(17), for the crossover.\\
6 & distinct symbols & $q$ is the carrier charge, $k$ the wavenumber, $h_{kij}$ the piezomagnetic constant, $b_n$ the mobility so that $\mu_{ik}$ stays permeability. Renamed in all three sets.\\
7 & order of the system & one sentence in each plan, at the end of the model section, handing the extra face conditions to the scholar. \\
\bottomrule
\end{tabular}
\end{center}

\section*{3. What the searches changed, and why the papers are narrower than they were}

This is the part that matters most, and it is not all good news.

\noindent\emph{Paper A.} Qu, Pan, Zhu, Jin and Roy published a thermoelectric theory of piezoelectric
semiconductors with the Seebeck and Peltier effects derived from a dissipation inequality
(\doi{10.1016/j.ijengsci.2022.103775}). Gupta and co-workers already put the modified gradient heat equation
into a piezo-semiconductor half space with voids and studied the Rayleigh wave
(\doi{10.1108/HFF-04-2025-0254}). Yang and Zhou already solved wave amplification in piezoelectric
semiconductor plates (\doi{10.1016/j.ijsolstr.2004.10.011}). So the claim that published PSC work closes at
most one direction of the carrier--temperature loop is false, and it is gone. What remains is the free-face
plate problem, the size effect and the attenuation curves.

\noindent\emph{Paper B.} A 2026 paper on generalized thermoelastic Lamb waves in one-dimensional hexagonal
quasicrystal nanoplates exists (\doi{10.1007/s10338-026-00755-4}). I could not read it. It is now week-one
reading for Neha, with four questions to answer before she writes a novelty sentence: which generalized
model, one lag or two, attenuation reported or not, phason dynamic or static. Until that is answered her
plan carries a provisional claim and says so.

\noindent\emph{Paper C.} Pyroelectric and pyromagnetic coefficients \emph{are} retained in plate work:
Kondaiah, Shankar and Ganesan (\doi{10.12989/csm.2013.2.1.001}) and Tocci Monaco and co-workers
(\doi{10.3390/nano11010087}) both keep them, the second inside a nonlocal strain gradient nanoplate. The
claim that they are rare is withdrawn, and so is the claim of a frequency-dependent magnetoelectric
coefficient, for which I found no source. What remains is the attenuation of Lamb modes in a layered plate
with a relaxation-time heat equation, with the effective magnetoelectric coefficient carried as a band from
two homogenisation routes. Matlam and Kedar (\doi{10.1080/01495739.2025.2612316}) and Kakar and Kakar
(\doi{10.12989/sem.2015.53.1201}) are the two nearest thermal neighbours and both are now cited.

\section*{4. What is not in the packages, and why}

There are no boundary conditions, no dispersion determinant, no solution method and no numerical value in
any of the six documents. That was your instruction and it is also the right division of labour: a scholar
who is handed a determinant learns nothing, and a scholar who is handed the model and the sequence of tests
has a job with a shape.

The internal folder carries the material that the packages deliberately lack:
\texttt{sumit\_equation\_source\_log.md}, \texttt{neha\_equation\_source\_log.md},
\texttt{anu\_equation\_source\_log.md}, the three benchmark files with every tolerance and every number we
hold, \texttt{changes\_v4\_review.md} and \texttt{style\_check.md}. The source logs mark each equation [R]
read, [S] secondary, [T] textbook, [D] your instruction, [P] proposed extension with no source, and they
name what I did not open. That list is short and uncomfortable: the Lim paper itself, the 2026 quasicrystal
thermoelastic paper, the Qu paper's equations, and the MGT heat equation as any of the cited authors write
it.

\section*{5. The week-one gates}

Each plan opens with reading, not computing, and each scholar reports on it before touching code. Sumit
confirms equations (17) to (19) of Lim et al.\ and the Peltier term in Qu et al. Neha obtains the 2026 paper
and writes down what is inside it. Anu reads the magnetoelectric sign convention in Chen, Pan and Chen and
the thermal face conditions in Kakar and Kakar. No calculations are set for the scholars yet, as agreed.

\noindent\textbf{Next step on my side, when you say so:} the three plans go out as the archives above, and
the first meeting is the reading report. The model document stays as the gate record.

\end{document}
